\HMTNarrativeHeading{De la récupération d’un état au transport d’une collection}

La récupération compatible permet d’étendre l’exposé depuis
un état jusqu’à des collections d’états. Une lecture enrichie
individualise ses antécédents lorsque sa mémoire distingue les
éléments qui partagent la même valeur visible. L’inverse
restitue alors l’état sur l’image de cette lecture.
La compatibilité avec le raffinement réunit les inverses finis
en une reconstruction d’histoires complètes.

Le transport des collections exige de conserver l’appartenance. Une
bijection d’histoires envoie chaque collection sur son image et chaque
élément sur son correspondant. Les unions, intersections et
compléments relatifs sont transportés au moyen de ces opérations.
La restriction de la bijection à une collection conserve son cardinal.
La théorie de la mémoire acquiert ainsi un lien précis avec les
opérations sur les ensembles d’histoires.

La mesure requiert en outre le transport des poids. Une symétrie compatible avec l'arbre de raffinement conserve la mesure des cylindres lorsqu'elle respecte les émissions et leurs multiplicités. La naturalité exprime la concordance entre transporter et tronquer. Les démonstrations suivantes explicitent ces relations sur les objets définis dans l'article.

L’étude spécifique de l’hypothèse du continu développe
la clôture généalogique et sa hiérarchie d’ensembles. Son résultat
cardinal est formulé sur cette clôture et son ensemble des parties interne.
La relation qu’incorpore cette partie est la récupération des
états et le transport des collections qui maintiennent la structure
du continu engendré. L’énoncé cardinal conserve son
développement spécialisé et son domaine expressément identifié.

Cette articulation situe le continu au sein de l’argument de
conservation. L’histoire complète, la collection à laquelle elle appartient
et la mesure de son cylindre requièrent des opérations liées,
aux contenus différents. L’article réunit leurs conditions de
compatibilité pour montrer comment la double projection se maintient
lors du passage d’une lecture individuelle à l’organisation des collections.
L’unité de l’exposé réside dans ces applications et dans les preuves
qui préservent leurs relations.
